\subsection{Definición de sustancia}

\section{Definición de mezcla}
    Para \cite{bib_termodinamica_cengel_2015}, una sustancia tiene composición química fija en cualquier parte del sistema, no tiene que estar necesariamente  conformada por un solo elemento químico o compuesto (sustancia). "Una mezcla de varias sustancias, también puede definirse como una sustancia, siempre que la mezcla sea homogénea". Una mezcla de dos o más fases de la misma sustancia, se considera como una sustancia pura, siempre que la composición química de las dos fases sea la misma; es asi que, una mezcla de hielo y agua líquida, es una sustancia pura porque ambas fases (sólida y líquida) tienen la misma composición química.

    qué es mezcla homogénea???

    De esta manera, se puede definirse el gas seco como una mezcla de sustancias (ejem. la mezcla de $O_2$, $N_2$ y otros gases en pequeñas cantidades) y considerarse como una sustancia homogénea

    \subsection{Mezcla de vapor de agua - aire}
        Para la mezcla binaria de vapor de agua y aire seco, 
        El aire seco es una mezcla de varios gases; pero, con frecuencia se considera como una sustancia pura, porque tiene una \textit{composición química uniforme} \parencite[p. 110]{bib_termodinamica_cengel_2015};
        
